\subsection{Measure Continuity and Regularity}

\begin{corollary} \label{cor:measure-continuity-up-down} \hlt{Measure Continuity}\\
Let $\mu$ be a measure, and let $(E_k)$ and $(F_k)$ be measurable sequences.
\begin{enumerate}[label=(\roman*)]
\setlength{\itemsep}{0pt}
\item If $E_k\uparrow E$, then $\mu(E_k)\uparrow\mu(E)$.
\item If $F_k\downarrow F$ and $\mu(F_1)<+\infty$, then $\mu(F_k)\downarrow\mu(F)$.
\end{enumerate}
\begin{proof}
\begin{enumerate}[label=(\roman*)]
\setlength{\itemsep}{0pt}
\item Set $G_1=E_1$ and $G_k=E_k\setminus E_{k-1}$ for $k\geq2$.\\
Then $E_k=\bigsqcup_{j=1}^{k}G_j$ and $E=\bigsqcup_{j\geq1}G_j$, so $\mu(E_k)=\sum_{j=1}^{k}\mu(G_j)\uparrow\sum_{j\geq1}\mu(G_j)=\mu(E)$.
\item Apply (i) to $H_k=F_1\setminus F_k\uparrow F_1\setminus F$. Since $\mu(F_1)<+\infty$, $\mu(F_k)=\mu(F_1)-\mu(H_k)\downarrow\mu(F_1)-\mu(F_1\setminus F)=\mu(F)$.
\end{enumerate}
\hlt{Structural Commentary}: Monotone set limits are disjoint series or complements of such series.\\
\hlt{Key Insight}: Downward continuity needs one finite anchor so subtraction is legal.
\end{proof}
\end{corollary}

\begin{definition} \label{def:set-liminf-limsup} \hlt{Limit Inferior and Superior of Sets}\\
For $(E_n)\subseteq\mathcal{P}(X)$, define $\limsup\limits_nE_n\defeq\bigcap\limits_{m=1}^{\infty}\bigcup\limits_{n\geq m}E_n$ and $\liminf\limits_nE_n\defeq\bigcup\limits_{m=1}^{\infty}\bigcap\limits_{n\geq m}E_n$.
\begin{enumerate}[label=(\roman*)]
\setlength{\itemsep}{0pt}
\item \hlt{Infinitely often}: $x\in\limsup\limits_nE_n$ iff $\{n:x\in E_n\}$ is infinite;
\item \hlt{Eventually always}: $x\in\liminf\limits_nE_n$ iff $\{n:x\notin E_n\}$ is finite;
\item \hlt{Set limit}: if $\limsup\limits_nE_n=\liminf\limits_nE_n$, write $\lim_nE_n$ for the common set.
\end{enumerate}
\end{definition}

\begin{proposition} \label{prop:set-limits-basic-laws} \hlt{Basic Laws for Set Limits}\\
Let $(E_n)\subseteq\mathcal{P}(X)$.
\begin{enumerate}[label=(\roman*)]
\setlength{\itemsep}{0pt}
\item $\emptyset\subseteq\liminf\limits_nE_n\subseteq\limsup\limits_nE_n\subseteq X$.
\item $E_n\uparrow E\Longrightarrow\liminf\limits_nE_n=\limsup\limits_nE_n=E=\bigcup\limits_nE_n$.
\item $E_n\downarrow E\Longrightarrow\liminf\limits_nE_n=\limsup\limits_nE_n=E=\bigcap\limits_nE_n$.
\item $(E_n)\subseteq\mathcal{M}\Longrightarrow\liminf\limits_nE_n,\limsup\limits_nE_n\in\mathcal{M}$.
\end{enumerate}
\begin{proof}
\begin{enumerate}[label=(\roman*)]
\setlength{\itemsep}{0pt}
\item Since $\bigcap\limits_m\bigcup\limits_{n\geq m}E_n\subseteq\bigcup\limits_{n\geq1}E_n\subseteq X$ and $\bigcap\limits_{n\geq m}E_n\subseteq\bigcup\limits_{n\geq m}E_n$ for all $m$, the inclusions follow.
\item If $E_n\uparrow E$, then $\{n:x\in E_n\}$ is nonempty iff it is cofinite, hence $\liminf\limits_nE_n=\limsup\limits_nE_n=\bigcup\limits_nE_n=E$.
\item If $E_n\downarrow E$, then $\{n:x\in E_n\}$ is infinite iff it is all of $\N$, hence $\liminf\limits_nE_n=\limsup\limits_nE_n=\bigcap\limits_nE_n=E$.
\item If $(E_n)\subseteq\mathcal{M}$, then $\bigcup\limits_{n\geq m}E_n,\bigcap\limits_{n\geq m}E_n\in\mathcal{M}$ $\forall m$; countable closure gives $\limsup\limits_nE_n,\liminf\limits_nE_n\in\mathcal{M}$.
\end{enumerate}
\hlt{Structural Commentary}: Set limsup means infinitely often; set liminf means eventually always.\\
\hlt{Key Insight}: Translate the definitions into pointwise membership statements.
\end{proof}
\end{proposition}

\begin{proposition} \label{prop:measure-liminf-limsup-inequalities} \hlt{Measure Bounds for Set Limits}\\
Let $(E_n)$ be measurable in a measure space $(X,\mathcal{M},\mu)$.
\begin{enumerate}[label=(\roman*)]
\setlength{\itemsep}{0pt}
\item $\mu(\liminf\limits_nE_n)\leq\liminf\limits_n\mu(E_n)$.
\item $\mu(\bigcup\limits_{n=1}^{\infty}E_n)<+\infty\Longrightarrow\limsup\limits_n\mu(E_n)\leq\mu(\limsup\limits_nE_n)$.
\end{enumerate}
\begin{proof}
\begin{enumerate}[label=(\roman*)]
\setlength{\itemsep}{0pt}
\item Set $A_m\defeq\bigcap\limits_{n\geq m}E_n$. Then $A_m\uparrow\liminf\limits_nE_n$ and $A_m\subseteq E_n$ for all $n\geq m$, so $\mu(A_m)\leq\inf\limits_{n\geq m}\mu(E_n)$.\\
By continuity from below, $\mu(\liminf\limits_nE_n)=\lim\limits_m\mu(A_m)\leq\sup\limits_m\inf\limits_{n\geq m}\mu(E_n)=\liminf\limits_n\mu(E_n)$.
\item Set $B_m\defeq\bigcup\limits_{n\geq m}E_n$. $B_m\downarrow\limsup\limits_nE_n$, $\mu(B_1)<+\infty$, $E_n\subseteq B_m$ $\forall n\geq m$, so $\sup\limits_{n\geq m}\mu(E_n)\leq\mu(B_m)$.\\
By continuity from above, $\limsup\limits_n\mu(E_n)=\inf\limits_m\sup\limits_{n\geq m}\mu(E_n)\leq\lim\limits_m\mu(B_m)=\mu(\limsup\limits_nE_n)$.
\end{enumerate}
\hlt{Structural Commentary}: Tail intersections increase to liminf; tail unions decrease to limsup.\\
\hlt{Key Insight}: Compare each tail envelope to the corresponding tail bounds of the measures.
\end{proof}
\end{proposition}

\begin{definition} \label{def:almost-everywhere-set-property} \hlt{Almost Everywhere}\\
Let $(X,\mathcal{M},\mu)$ be a measure space, let $E\in\mathcal{M}$, and let $P(x)$ be a property on $E$.
\begin{enumerate}[label=(\roman*)]
\setlength{\itemsep}{0pt}
\item $P$ holds \hlt{$\mu$-almost everywhere on $E$} iff $\exists N\in\mathcal{M}$ with $N\subseteq E$, $\mu(N)=0$, and $P$ holds on $E\setminus N$.
\item For Lebesgue measure, write \hlt{almost everywhere}.
\end{enumerate}
\end{definition}

\begin{theorem} \label{thm:borel-cantelli-summable-events} \hlt{First Borel--Cantelli Lemma}\\
Let $(X,\mathcal{M},\mu)$ be measure space, $(E_n)\subseteq\mathcal{M}$, set $L\defeq\limsup\limits_nE_n=\bigcap\limits_{m\geq1}\bigcup\limits_{n\geq m}E_n$. Assume $\sum\limits_{n=1}^{\infty}\mu(E_n)<+\infty$.
\begin{enumerate}[label=(\roman*)]
\setlength{\itemsep}{0pt}
\item $L\in\mathcal{M}$ and $\mu(L)=0$.
\item $\mu$-almost every $x\in X$ belongs to at most finitely many $E_n$.
\end{enumerate}
\begin{proof}
\begin{enumerate}[label=(\roman*)]
\setlength{\itemsep}{0pt}
\item Set $B_m\defeq\bigcup\limits_{n\geq m}E_n$. Then $B_m\downarrow L$ and $\mu(B_m)\leq\sum\limits_{n\geq m}\mu(E_n)$.\\
Since $\mu(B_1)\leq\sum\limits_{n\geq1}\mu(E_n)<+\infty$, continuity from above gives $\mu(L)=\lim\limits_{m\to\infty}\mu(B_m)\leq\lim\limits_{m\to\infty}\sum\limits_{n\geq m}\mu(E_n)=0$.
\item Since $x\in L\iff\forall m\,\exists n\geq m$ with $x\in E_n$, $L$ is the set of points lying in infinitely many $E_n$.\\
By (i), outside the null set $L$, membership occurs only finitely often.
\end{enumerate}
\hlt{Structural Commentary}: Summable tails force the limsup set to have zero measure.\\
\hlt{Key Insight}: Points appearing infinitely often must lie in every tail union.
\end{proof}
\end{theorem}

\begin{remark} \label{rem:borel-null-regular-measurable-sets} \hlt{Borel Sets, Null Sets, and Regular Approximation}\\
Open, closed, and Borel subsets of $\R^p$ are measurable. Each $A\in\mathcal{M}(\mu)$ admits closed--open approximations $F\subseteq A\subseteq G$ with $\mu(G\setminus A)$ and $\mu(A\setminus F)$ arbitrarily small. In Lebesgue measure, one may take $F_\sigma\subseteq A\subseteq G_\delta$ with $\mu(G_\delta\setminus F_\sigma)=0$. Null sets form a $\sigma$-ring, and countable subsets of $\R^p$ are null.\\
\hlt{Structural Commentary}: Measurable sets are Borel sets modified by null sets.\\
\hlt{Key Insight}: Regularity explains why open and closed approximations are enough for integration.
\end{remark}

\begin{proposition} \label{prop:local-measurability-bounded-cubes} \hlt{Local Measurability on Bounded Cubes}\\
For arbitrary $E\subseteq\R^p$, $E\in\mathcal{L}^p\Longleftrightarrow E\cap[-n,n]^p\in\mathcal{L}^p$ for all $n\in\N$.
\begin{proof}
\begin{enumerate}[label=\roman*.]
\setlength{\itemsep}{0pt}
\item[$(\Rightarrow)$] If $E\in\mathcal{L}^p$, then $E\cap[-n,n]^p\in\mathcal{L}^p$ for all $n$, since $[-n,n]^p$ is Borel.
\item[$(\Leftarrow)$] If $E\cap[-n,n]^p\in\mathcal{L}^p$ for all $n$, then $E=\bigcup_{n=1}^{\infty}(E\cap[-n,n]^p)\in\mathcal{L}^p$.
\end{enumerate}
\hlt{Structural Commentary}: Nested Borel cubes exhaust $\R^p$.\\
\hlt{Key Insight}: Measurability is local along a countable exhaustion.
\end{proof}
\end{proposition}

\begin{proposition} \label{prop:real-line-local-measurability-unit-intervals} \hlt{Local Measurability on Unit Intervals}\\
Let $I_n\defeq(n,n+1]$, $n\in\Z$. For $E\subseteq\R$, $E\in\mathcal{L}^1$ iff $E\cap I_n\in\mathcal{L}^1$ for all $n\in\Z$.
\begin{proof}
\begin{enumerate}[label=\roman*.]
\setlength{\itemsep}{0pt}
\item[$(\Rightarrow)$] Since $I_n\in\mathcal{B}(\R)\subseteq\mathcal{L}^1$, $E\in\mathcal{L}^1$ implies $E\cap I_n\in\mathcal{L}^1$ for all $n$.
\item[$(\Leftarrow)$] If $E\cap I_n\in\mathcal{L}^1$ for all $n$, then $E=\bigcup_{n\in\Z}(E\cap I_n)\in\mathcal{L}^1$.
\end{enumerate}
\hlt{Structural Commentary}: A countable Borel exhaustion localizes measurability.\\
\hlt{Key Insight}: Local measurable pieces patch by countable unions.
\end{proof}
\end{proposition}

\begin{proposition} \label{prop:finite-measure-meets-every-unit-interval} \hlt{Finite Measure Meeting Every Unit Interval}\\
With $I_n\defeq(n,n+1]$, there exists $E\in\mathcal{B}(\R)$ such that $0<m(E\cap I_n)$ for all $n\in\Z$ and $m(E)<+\infty$.
\begin{proof}
Set $\varepsilon_n\defeq2^{-\abs{n}-3}$ and $E\defeq\bigcup_{n\in\Z}(n,n+\varepsilon_n]$. Then $(n,n+\varepsilon_n]\subseteq I_n$ are pairwise disjoint Borel intervals, so $E\in\mathcal{B}(\R)$ and $m(E\cap I_n)=\varepsilon_n>0$. Countable additivity gives $m(E)=\sum_{n\in\Z}\varepsilon_n=\frac38<+\infty$.\\
\hlt{Structural Commentary}: Finite measure can be spread over countably many cells.\\
\hlt{Key Insight}: Positive local masses can be summable globally.
\end{proof}
\end{proposition}

\begin{proposition} \label{prop:null-sets-sigma-ideal} \hlt{Null Sets Form a Sigma-Ideal}\\
Let $(X,\mathcal{M},\mu)$ be a measure space and $\mathcal{Z}\defeq\{N\in\mathcal{M}:\mu(N)=0\}$.\\
Then $\mathcal{Z}$ is a $\sigma$-ring, and $N\in\mathcal{Z}$, $E\in\mathcal{M}$ imply $N\cap E\in\mathcal{Z}$.
\begin{proof}
If $(N_k)\subseteq\mathcal{Z}$, then $\mu(\bigcup_kN_k)\leq\sum_k\mu(N_k)=0$, so $\bigcup_kN_k\in\mathcal{Z}$.\\
If $A,B\in\mathcal{Z}$, then $A\setminus B\subseteq A$, hence $A\setminus B\in\mathcal{Z}$. Finally, $N\cap E\subseteq N$ gives $N\cap E\in\mathcal{Z}$.\\
\hlt{Structural Commentary}: Null sets are stable under countable unions and measurable cuts.\\
\hlt{Key Insight}: Monotonicity makes every measurable piece of a null set null.
\end{proof}
\end{proposition}

\begin{proposition} \label{prop:completed-sigma-algebra} \hlt{Completed Sigma-Algebra}\\
Let $(X,\mathcal{M},\mu)$ be measure space, $\mathcal{Z}\defeq\{N\in\mathcal{M}:\mu(N)=0\}$; $\overline{\mathcal{M}}_\mu\defeq\{A\subseteq X:\exists E\in\mathcal{M},\ \exists N\in\mathcal{Z},\ A\triangle E\subseteq N\}$.
\begin{enumerate}[label=(\roman*)]
\setlength{\itemsep}{0pt}
\item $\mathcal{M}\subseteq\overline{\mathcal{M}}_\mu$, and $\overline{\mathcal{M}}_\mu$ is a $\sigma$-algebra;
\item $A\subseteq N\in\mathcal{Z}\Longrightarrow A\in\overline{\mathcal{M}}_\mu$;
\item $\overline{\mathcal{M}}_\mu=\{E\cup Z:E\in\mathcal{M},\ \exists N\in\mathcal{Z},\ Z\subseteq N\}$.
\end{enumerate}
\begin{proof}
\begin{enumerate}[label=(\roman*)]
\setlength{\itemsep}{0pt}
\item For $E\in\mathcal{M}$, $E\triangle E=\emptyset$, so $\mathcal{M}\subseteq\overline{\mathcal{M}}_\mu$. If $A\triangle E\subseteq N$, then $A^c\triangle E^c=A\triangle E\subseteq N$.\\
If $A_k\triangle E_k\subseteq N_k$, then $(\bigcup_kA_k)\triangle(\bigcup_kE_k)\subseteq\bigcup_kN_k\in\mathcal{Z}$. Hence $\overline{\mathcal{M}}_\mu$ is a $\sigma$-algebra.
\item If $A\subseteq N\in\mathcal{Z}$, then $A\triangle\emptyset\subseteq N$, so $A\in\overline{\mathcal{M}}_\mu$.
\item If $A\triangle E\subseteq N$, then $A=(E\setminus N)\cup(A\cap N)$, with $E\setminus N\in\mathcal{M}$ and $A\cap N\subseteq N$.\\
Conversely, if $A=E\cup Z$ and $Z\subseteq N\in\mathcal{Z}$, then $A\triangle E\subseteq Z\subseteq N$.
\end{enumerate}
\hlt{Structural Commentary}: Completion adds all null modifications of measurable sets.\\
\hlt{Key Insight}: Equality modulo null sets is the membership criterion.
\end{proof}
\end{proposition}

\begin{proposition} \label{prop:null-ideal-and-measure-completion} \hlt{Completed Measure}\\
Let $(X,\mathcal{M},\mu)$ be a measure space, $\mathcal{Z}\defeq\{N\in\mathcal{M}:\mu(N)=0\}$,\\
and $\overline{\mathcal{M}}_\mu\defeq\{A\subseteq X:\exists E\in\mathcal{M},\ \exists N\in\mathcal{Z},\ A\triangle E\subseteq N\}$. For $A\in\overline{\mathcal{M}}_\mu$, define $\overline{\mu}(A)\defeq\mu(E)$ whenever $E\in\mathcal{M}$ and $A\triangle E\subseteq N\in\mathcal{Z}$. Then $\overline{\mu}$ is a well-defined complete measure on $\overline{\mathcal{M}}_\mu$, and $\overline{\mu}|_{\mathcal{M}}=\mu$.
\begin{proof}
If $A\triangle E\subseteq N_1$ and $A\triangle F\subseteq N_2$, with $E,F\in\mathcal{M}$ and $N_1,N_2\in\mathcal{Z}$, then $E\triangle F\subseteq N_1\cup N_2$, so $\mu(E)=\mu(F)$. Hence $\overline{\mu}$ is well defined; since $E\triangle E=\emptyset$, $\overline{\mu}|_{\mathcal{M}}=\mu$.\\
For disjoint $(A_k)\subseteq\overline{\mathcal{M}}_\mu$, choose $E_k\in\mathcal{M}$ and $N_k\in\mathcal{Z}$ with $A_k\triangle E_k\subseteq N_k$, and set $N\defeq\bigcup_kN_k$.\\
Then $E_i\cap E_j\subseteq N$ for $i\neq j$. Put $F_1\defeq E_1$ and $F_k\defeq E_k\setminus\bigcup_{j<k}E_j$.\\
Then $(F_k)$ disjoint, $A_k\triangle F_k\subseteq N$, $(\bigcup_kA_k)\triangle(\bigcup_kF_k)\subseteq N$, so $\overline{\mu}\left(\bigcup_kA_k\right)=\mu\left(\bigcup_kF_k\right)=\sum_k\mu(F_k)=\sum_k\overline{\mu}(A_k)$.\\
If $A\subseteq B\in\overline{\mathcal{M}}_\mu$ and $\overline{\mu}(B)=0$, choose $E\in\mathcal{M}$ and $N\in\mathcal{Z}$ with $B\triangle E\subseteq N$.\\
Then $\mu(E)=0$, so $A\subseteq B\subseteq E\cup N\in\mathcal{Z}$; hence $A\in\overline{\mathcal{M}}_\mu$ and $\overline{\mu}(A)=0$.\\
\hlt{Structural Commentary}: The completed measure ignores null errors.\\
\hlt{Key Insight}: Null perturbations do not change countable additivity.
\end{proof}
\end{proposition}
